\section{Introduction}\label{sec:intro}

A bank that deploys a machine-learning model into a decision-relevant role is, in the United States since 2011 and in the European Union since 2024, expected to maintain a record of every prediction the model produces: what code computed it, on what data, with what hyperparameters, in what compute environment, with what random seed, at what time. The expectations are codified in Federal Reserve SR 11-7 / OCC Bulletin 2011-12 \citep{FedSR117, OCC2011} and in the EU AI Act \citep{EUAIAct}, Articles 12 and 19. The compliance question is binary: either the audit trail exists, is intact, and is replayable, or it does not.

The technical question is more subtle. The compliance literature describes \emph{what} should be in an audit trail; it says little about \emph{how} to actually log one for a Python ML pipeline at scale. Existing open-source tools fall into three categories. (1) Experiment-tracking platforms (MLflow, Weights \& Biases) are designed for the train-time exploration phase, with rich UIs for comparing runs but limited support for inference-time logging at scale and no built-in immutability or hash-chained integrity. (2) Data and model versioning tools (DVC, Pachyderm) version artifacts but do not log per-prediction inputs and outputs. (3) Cloud-managed monitoring services (AWS SageMaker Model Monitor, Google Vertex AI Model Monitoring) are vendor-locked, expensive at audit-grade volume, and require trusting the cloud provider with the exact inference data the audit is designed to protect.

What is missing is an open-source, vendor-neutral Python package that produces an SR 11-7-shaped audit trail with low-friction integration. We provide one and demonstrate it.

\subsection{Contribution}

This paper makes two contributions.

First, we introduce \texttt{mr-audit}, an open-source Python package (MIT-licensed; PyPI \texttt{mr-audit}; repository \href{https://github.com/ayk5511/mr-audit}{github.com/ayk5511/mr-audit}) that produces immutable, hash-chained, replayable, drift-aware audit trails for ML inference pipelines. The package is integrated into a calling pipeline by a one-line decorator (\verb|@auditable|) and a context manager (\texttt{AuditContext}); no central server, no cloud account, no schema-design effort. Records are stored locally in either SQLite or JSON Lines format, both vendor-neutral. The package implements a SHA-256 hash chain (each record's identifier is the canonical-JSON SHA-256 of its content; each record carries the previous record's identifier; tampering with any single record breaks the chain). It ships with a replay engine that reconstructs any prediction from the log, a drift-detection module (Kolmogorov-Smirnov two-sample test plus Population Stability Index with Laplace smoothing), and a regulator-export bundle that produces a self-contained zip archive whose integrity can be verified using off-the-shelf hash tools without any \texttt{mr-audit} code installed.

Second, we demonstrate the package on a real ML pipeline by retrospectively instrumenting the seven-model volatility-forecasting panel of \citet{KhanVol2026} (S\&P 500 5-day realized volatility forecasts from three GARCH variants, HAR-RV, LightGBM, XGBoost, and an equal-weight ensemble across 980 trading days, 2022-01-03 to 2025-11-26). The same panel is the empirical demonstration of \citet{Khan2026VolEval}, which uses \texttt{vol-eval} to apply Diebold-Mariano, Model Confidence Set, and Hansen SPA tests to the seven models. \texttt{mr-audit} and \texttt{vol-eval} are intended as complementary infrastructure: \texttt{vol-eval} addresses the question ``did your forecasting comparison use the right significance test?''; \texttt{mr-audit} addresses ``can you prove tomorrow that today's prediction was made with this code, this data, and these parameters?''. The instrumented run produces a 6,825-record audit log. We then exercise the five primary use cases an SR 11-7 model risk management team would run on such a log:

\begin{enumerate}[leftmargin=*]
    \item \textbf{Reproducibility audit.} A randomly-sampled set of 5 records replays bit-identically: the original input hash, the model parameters, and the final prediction match exactly when re-run.
    \item \textbf{Drift detection.} The KS test flags drift in 44 of the 44 monthly windows evaluated against the 2022 Q1 baseline. The pattern is consistent with the documented regime shift between high-vol 2022 and lower-vol 2023+ in \citet{KhanVol2026}.
    \item \textbf{Validation report.} Per-model summary statistics (n, prediction mean, prediction std, library versions at last record) are computable from the audit log alone, with no external data.
    \item \textbf{Regulator export.} A 841 KB zip bundle is produced from the 6,825-record log; its manifest, file SHA-256 hashes, and chain integrity all verify cleanly using a stand-alone tool that imports nothing from \texttt{mr-audit}.
    \item \textbf{Right-to-explanation.} Any historical prediction can be reconstructed from the log alone into a complete provenance bundle (code commit, library versions, model parameters, input hash, and value), satisfying GDPR Article 22 and EU AI Act Article 86 explanation requirements with hash-chain-verifiable integrity.
\end{enumerate}

\subsection{Related work}\label{sec:related_work}

This paper sits at the intersection of three literatures: ML model documentation, internal algorithmic auditing, and supervised model risk management. We engage with each.

\textit{Model documentation.} \citet{MitchellModelCards2019} introduce Model Cards as a structured documentation format for trained ML models, capturing intended uses, performance metrics, and ethical considerations. Subsequent work extends the pattern to datasets \citep{GebruDatasheets2021, PushkarnaDataCards2022} and to industry-supplier declarations \citep{ArnoldFactSheets2019}. These are \emph{point-in-time} documentation artefacts produced once per model. They are necessary but not sufficient for SR 11-7 / EU AI Act compliance, which require \emph{per-prediction} provenance maintained across the deployed lifetime of the system. Our contribution complements model documentation: where Model Cards describe what the model is, the audit trail described here records what each prediction was.

\textit{Algorithmic auditing.} \citet{Raji2020Audit} propose an end-to-end internal-auditing framework spanning scoping, mapping, artefact collection, testing, and reflection. They emphasise that compliance-grade auditing must produce verifiable artefacts at each stage of the model lifecycle, not just at deployment. \citet{BrundageVerifiableClaims2020} argue that ``verifiable claims'' about AI system behaviour require mechanisms beyond self-reporting: third-party audits, bias and safety bounties, sharing of incident information, and tamper-evident audit trails are all listed as concrete mechanisms. Our work operationalises the audit-trail mechanism specifically, in a form designed against the existing record-keeping clauses of SR 11-7 and EU AI Act Article 19. \citet{Sambasivan2021Cascades} document the practitioner-side observation that data and provenance work is systematically under-resourced relative to modelling work; we take this as motivating the need for low-friction (one-decorator) integration.

\textit{MLOps and model management.} \citet{SchelterContinuousIntegration2018} survey the engineering challenges of ML model management, identifying audit trails as one of seven open problems. The intervening seven years have produced extensive tooling for the other six (versioning, experiment tracking, serving, monitoring, feature stores, model registries) but the audit-trail layer specifically remains under-served, particularly for the regulated-industry use case where the trail must be cryptographically verifiable and regulator-portable. The capability comparison in \Cref{tab:tool_comparison} substantiates this gap quantitatively.

The narrow positioning of this paper is: a Python package and supporting schema that satisfies the per-prediction audit-trail requirements of SR 11-7, EU AI Act Articles 12 and 19, GDPR Article 22, and NIST AI 100-1, with integration cost low enough that the cost-of-not-doing-it (regulatory exposure plus the data-cascade penalty of \citealp{Sambasivan2021Cascades}) is the dominant economic argument.

\subsection{What this paper claims and does not claim}

Three explicit scoping decisions deserve flagging.

First, we do \emph{not} claim that using \texttt{mr-audit} makes any model \emph{compliant} with SR 11-7, OCC 2011-12, or the EU AI Act. Compliance is a determination by the firm and its supervisor; the package provides the technical record-keeping infrastructure that compliance requires.

Second, the empirical demonstration is a \emph{retrospective} instrumentation. The 6,825 records are written today against the previously-fit forecast panel of \citet{KhanVol2026}; we are not running the seven-model pipeline live and recording in real time. The retrospective form is appropriate for the demonstration's purpose (showing that the five use cases work end-to-end at a realistic scale of 6{,}825 records) but does not establish that \texttt{mr-audit} would handle a live high-throughput inference workload without further engineering. Section~\ref{sec:limitations} discusses this honestly.

Third, the EU AI Act's high-risk classification (Annex III) does not currently include volatility forecasting. SR 11-7 is the live regulatory pull for the demonstration's empirical target; the EU AI Act framing is forward-looking. We do not overclaim that vol forecasting is presently in scope of the AI Act's high-risk regime.

\subsection{Roadmap}

\Cref{sec:regulatory} reviews the SR 11-7 and EU AI Act requirements relevant to audit trails, distinguishing what the regulations specify from what they leave to the firm. \Cref{sec:package} describes the \texttt{mr-audit} package: design principles, schema, hash chain, storage backends, replay engine, drift module, export bundle, and CLI. \Cref{sec:demonstration} reports the empirical demonstration on the Khan (2026) volatility panel, including the five use cases and the scaling-throughput results. \Cref{sec:discussion} discusses the implications and limitations. \Cref{sec:reproducibility} states the reproducibility commitment, including the audit script that re-derives every numerical claim in this paper.
